% Answers to Chapter 10, from lectures/L10/appendix/c_solutions.md; the self-test from
% lectures/L10/README.md.

\facitavsnitt{c10}{Kapitel 10}{Exponentialfunktioner, logaritmer och decibel}
\begin{facit}

\svar{1.1} \sv{a} $\frac{\pi}{12}$ \sv{b} $-\frac{5\pi}{12}$

\svar{1.2} \sv{a} $-30^{\circ}$ \sv{b} $270^{\circ}$

\svar{1.3} $u(t) = 4\sin\left(50\pi t - \frac{\pi}{4}\right)\,\text{V}$

\svar{1.4} $u(t) = 4\sin\left(100\pi t - \frac{\pi}{6}\right)\,\text{V}$. Kurvan har
perioden $T = 20\,\text{ms}$ och börjar i $u(0) = -2\,\text{V}$. Den går uppåt genom noll vid
$t \approx 1{,}67\,\text{ms}$ (en tolftedels period för sent), når toppen $4\,\text{V}$ vid
$6{,}67\,\text{ms}$ och botten $-4\,\text{V}$ vid $16{,}67\,\text{ms}$.

\svar{1.5} $\delta \approx -3{,}25\,\text{rad}$ eller $\delta \approx -1{,}15\,\text{rad}$,
exakt $-\frac{31\pi}{30}$ och $-\frac{11\pi}{30}$.

\svar{2.1} \sv{a} $3$ \sv{b} $10$ \sv{c} $\approx 3{,}79$

\svar{2.2} $t \approx 132{,}88\,\text{h}$, ungefär $132$~timmar och $53$~minuter.

\svar{2.3} $G_{\text{lin}} \approx 50{,}12$, alltså cirka $50$.

\svar{2.4} $\abs{U} \approx 10{,}0\,\text{V}$ ($U_{\text{RMS}} \approx 7{,}08\,\text{V}$).

\svar{3.1} \sv{a} $2$ \sv{b} $3$ \sv{c} $0$ \sv{d} $-2$ \sv{e} $1$ \sv{f} $0$

\svar{3.2} \sv{a} $\log 12$ \sv{b} $\log 5$ \sv{c} $\log 25$ \sv{d} $1$ \sv{e} $\log 2$

\svar{3.3} \sv{a} $5$ \sv{b} $-3$ \sv{c} $\approx 2{,}86$ \sv{d} $\approx 3{,}00$ \sv{e} $2$

\svar{3.4} \sv{a} $100$ \sv{b} $1$ \sv{c} $5$ \sv{d} $100$
\sv{e} $\frac{e}{2} \approx 1{,}36$

\svar{3.5} \sv{a} $40\,\text{dB}$ \sv{b} $0\,\text{dB}$ \sv{c} $-20\,\text{dB}$
\sv{d} Dämpning: utsignalen är mindre än insignalen, så kvoten är mindre än $1$ och dess
logaritm negativ.

\svar{3.6} \sv{a} $10$ \sv{b} $1000$ \sv{c} $\approx 2{,}00$ \sv{d} $\approx 0{,}708$
\sv{e} $\approx 1{,}0\,\text{V}$

\svar{3.7} \sv{a} $20\,\text{dB}$ \sv{b} $\approx -3{,}01\,\text{dB}$
\sv{c} Effekten är proportionell mot spänningen i kvadrat, $P = U^2/R$, och exponenten
$2$ blir en faktor framför logaritmen: $10 \cdot 2 = 20$.
\sv{d} Effekten: $\approx 3\,\text{dB}$ ($3{,}01$). Spänningen: $\approx 6\,\text{dB}$ ($6{,}02$).

\svar{3.8} \sv{a} $26\,\text{dB}$ \sv{b} $\approx 20{,}0$
\sv{c} $3{,}98 \times 7{,}94 \times 0{,}631 \approx 20{,}0$
\sv{d} Linjära förstärkningar multipliceras men dB-värden adderas, eftersom
$\log(ab) = \log a + \log b$.

\svar{3.9} \sv{a} $0\,\text{dBV}$ \sv{b} $20\,\text{dBV}$ \sv{c} $-20\,\text{dBV}$
\sv{d} $\approx 3{,}98\,\text{V}$ \sv{e} $\approx 5{,}63\,\text{V}$

\svar{3.10} \sv{a} $a = 0{,}98$ \sv{b} $\approx 34{,}3\,\text{h}$ \sv{c} $\approx 7{,}3$ dygn
\sv{d} $C \cdot a^t = 2C$ ger $a^t = 2$, och logaritmering ger $t\log a = \log 2$.

\svar{3.11} \sv{a} Det är $\log(ab)$ som är $\log a + \log b$; $\log(a + b)$ kan inte förenklas.
\sv{b} Kvoten är ett basbyte: $\frac{\log 8}{\log 2} = \log_2 8 = 3$.
\sv{c} Exponenten blir en faktor: $\log(a^n) = n\log a$.
\sv{d} $0\,\text{dB}$ ger $G_{\text{lin}} = 1$: signalen passerar oförändrad.
\sv{e} Tiopotens och tiologaritm tar ut varandra: $10^{\log x} = x$.
\sv{f} $\ln 0$ är odefinierat; det är $\ln 1$ som är $0$.

\svar[Testa dig själv]{T} \sv{1} $x = 3$ \sv{2} $G_{\text{lin}} = 100$

\end{facit}
